\documentclass{article}

% NeurIPS-style packages
% \usepackage[final]{neurips_2024}   % use neurips_2024.sty if available; else article
\usepackage{amsmath, amssymb}
\usepackage{graphicx}
\usepackage{hyperref}
\usepackage{url}
\usepackage{booktabs}
\usepackage{microtype}
\usepackage{xcolor}
\usepackage{algorithm}
\usepackage{algorithmic}

\hypersetup{colorlinks=true, linkcolor=blue!70!black, citecolor=blue!70!black}

\graphicspath{{figures/}}

% If neurips_2024.sty is not available, comment the \usepackage above and
% uncomment the following minimal style:
%\usepackage[margin=2.5cm]{geometry}

% ----------------------------------------------------------------
\title{PhotoMedGemma: Machine-Precision Compilation of a\\
       Medical Vision-Language Model onto Photonic Hardware}

\author{%
    Dalton Omondi\\[2pt]
  \small Electrical and Electronics Engineering, Kenyatta University\\[2pt]
  \small \href{https://github.com/GabuGravin41/photonmedgemma}{github.com/GabuGravin41/photonmedgemma}
  % \thanks{Independent Research. E-mail: photomedgemma@github.com}
}

% ----------------------------------------------------------------
\begin{document}
\maketitle

% ----------------------------------------------------------------
\begin{abstract}
We present PhotoMedGemma, a simulation-validated open-source compiler that
maps the weight matrices of MedGemma-4B-IT --- a 4-billion-parameter multimodal
medical AI --- to phase settings for silicon photonic Mach-Zehnder
interferometer (MZI) meshes.
Photonic chips perform matrix-vector multiplication using light, projecting
$\sim 37\times$ better energy efficiency than GPU inference under idealised
steady-state assumptions.
The critical barrier to deployment of AI in global-south hospitals
(data sovereignty, connectivity, energy) motivates this hardware approach.

Our compiler performs truncated singular value decomposition (SVD) of each
weight matrix, followed by Clements decomposition of the resulting unitary
factors into MZI phase angles.
We validate the pipeline numerically on real MedGemma activations captured
from five breast cancer PNG images.
The compiled Clements mesh reproduces both unitary factors to errors
$< 8 \times 10^{-15}$ (machine precision) for all four layer-0 attention
projections (query, key, value, output).
We report a rank sweep from $r = 64$ to full rank, showing
SVD approximation error decreasing from $6.9 \times 10^{-1}$ at rank 64
to machine precision at full rank, with MZI counts scaling as $r(r-1)/2$.
All code and compiled phase maps are released publicly.
\end{abstract}

% ----------------------------------------------------------------
\section{Introduction}

Large language models trained for medicine, such as MedGemma~\citep{medgemma2024},
achieve clinical-grade accuracy on diagnostic tasks.
Yet their deployment in low-resource settings --- rural hospitals, field
clinics, regional health systems in Africa and South Asia --- is blocked by
energy requirements ($\sim 300$\,W/GPU), cloud API costs, and data sovereignty
constraints.

Photonic neural networks~\citep{shen2017, bandyopadhyay2022} offer a
fundamentally different compute substrate: matrix-vector multiplication encoded
into the interference pattern of light in an MZI mesh.
For a fixed, trained model, all weight matrices are constant.
The key insight of this work is that \textbf{static photonic compilation} ---
encoding weights permanently into MZI phase settings --- converts inference
into a passive physical process: light in, computation by geometry,
light out.
No dynamic weight updates are needed; the chip's heaters are set once and
consume near-zero steady-state power.

\paragraph{Contributions.}
The SVD-Clements pipeline for photonic neural networks is well-established
\citep{reck1994, clements2016, shen2017, pai2023}.
Our contribution is its application to a \emph{real production-scale medical
VLM}, with:
\begin{itemize}
  \item[(1)] A simulation-validated open-source compiler from MedGemma weights
             to Clements MZI phase maps, verified at machine precision
             ($\varepsilon < 8 \times 10^{-15}$).
  \item[(2)] A simulation framework applying compiled phases to real MedGemma
             activations from breast cancer images, measuring SVD approximation
             error against digital ground truth.
  \item[(3)] A rank sweep ($r = 64, 128, 256, 512, 1024, 2048$) with
             honest reporting of accuracy-hardware trade-offs.
  \item[(4)] Open-source release of all code, compiled phases, and data.
\end{itemize}

% ----------------------------------------------------------------
\section{Background}

\subsection{MZI Mesh as a Unitary Operator}

A passive linear optical network on $N$ modes implements a unitary matrix
$U \in \mathcal{U}(N)$ on the complex mode amplitudes~\citep{reck1994}.
The \textbf{Clements decomposition}~\citep{clements2016} factors any
$U \in \mathcal{U}(N)$ into exactly $N(N-1)/2$ two-mode rotations:
\begin{equation}
  U = D \cdot T_1 \cdot T_2 \cdots T_{N(N-1)/2},
  \quad
  T(\theta,\varphi) = \begin{pmatrix}
    \cos\tfrac{\theta}{2} & -e^{-i\varphi}\sin\tfrac{\theta}{2} \\
    e^{i\varphi}\sin\tfrac{\theta}{2} & \cos\tfrac{\theta}{2}
  \end{pmatrix},
  \label{eq:clements}
\end{equation}
where $D$ is a diagonal phase screen.
The forward pass applies MZIs in descending column order, then $D$.

\subsection{SVD-Based Photonic Weight Compilation}

A weight matrix $W \in \mathbb{R}^{m\times n}$ is compiled via SVD:
$W = U\Sigma V^\dagger$.
$V^\dagger$ and $U$ are unitary and implemented as Clements meshes;
$\Sigma = \text{diag}(\sigma_1,\ldots,\sigma_k)$ is implemented as a
variable optical attenuator (VOA) bank.
For a rank-$r$ approximation, we retain the top-$r$ singular triplets.

% ----------------------------------------------------------------
\section{The PhotoMedGemma Compiler}
\label{sec:compiler}

\subsection{Pipeline}

\begin{algorithm}[H]
\caption{PhotoMedGemma Compilation}
\label{alg:compiler}
\begin{algorithmic}[1]
  \STATE \textbf{Input}: weight matrix $W \in \mathbb{R}^{m\times n}$, rank $r$
  \STATE Compute truncated SVD: $U_r, \Sigma_r, V_r^\dagger \leftarrow \text{SVD}(W, r)$
  \STATE Embed $V_r^\dagger$ in a full $n\times n$ unitary $\hat{V}^\dagger$
  \STATE Apply Clements left-elimination to $\hat{V}^\dagger$:
         $\{(\theta_k^V, \varphi_k^V)\}_{k=1}^{n(n-1)/2}$, $\boldsymbol{\psi}^V$
  \STATE Embed $U_r$ in a full $m\times m$ unitary $\hat{U}$
  \STATE Apply Clements left-elimination to $\hat{U}$:
         $\{(\theta_k^U, \varphi_k^U)\}_{k=1}^{m(m-1)/2}$, $\boldsymbol{\psi}^U$
  \STATE Quantise all phases to $B=12$-bit DAC codes
  \STATE Write phase map to \texttt{phase\_map.json}; GDS to \texttt{.gds}
  \STATE \textbf{Output}: phase settings for $\hat{V}^\dagger$ mesh,
         attenuators $\sigma_1,\ldots,\sigma_r$, phase settings for $\hat{U}$ mesh
\end{algorithmic}
\end{algorithm}

\subsection{Clements Convention}

The correct transfer matrix for the left-elimination Clements algorithm is
equation~(\ref{eq:clements}), applied in \textbf{descending column order}
during the forward pass.
This is verified by requiring $U_{\text{reconstructed}}\mathbf{x} = U_{\text{target}}\mathbf{x}$
to machine precision before writing the phase map.

% ----------------------------------------------------------------
\section{Experiments}

\subsection{Setup}

We ran MedGemma-4B-IT on 5 breast cancer PNG images using a Kaggle T4 GPU
(4-bit NF4 quantisation).
For each image, we captured layer-0 attention projection inputs and outputs,
plus weight matrices $W_q, W_k, W_v, W_o$ (dequantised from NF4 to float32).
Three metrics are computed:

\noindent \textbf{Chip fidelity:}
$\varepsilon_{\rm chip} = \|\text{MZI\_sim}(\mathbf{x}) - U_r V_r^\dagger\mathbf{x}\|
/ \|U_r V_r^\dagger\mathbf{x}\|$.

\noindent \textbf{SVD approximation error:}
$\varepsilon_{\rm SVD}(r) = \|W_r\mathbf{x} - W\mathbf{x}\| / \|W\mathbf{x}\|$.

\noindent \textbf{Quantisation error:}
$\varepsilon_{\rm quant} = \|W_{\rm kaggle}\mathbf{x} - \mathbf{y}_{\rm MG}\|
/ \|\mathbf{y}_{\rm MG}\|$.

\subsection{Compilation Fidelity}

\begin{table}[h]
\centering
\caption{Compilation fidelity for all four layer-0 attention projections.
         All errors at machine precision ($\approx 10^{-15}$).}
\label{tab:fidelity}
\begin{tabular}{lccccc}
\toprule
\textbf{Proj.} & \textbf{Shape} & \textbf{Chip fid.} & \textbf{SVD err.\ ($r$=64)} &
\textbf{SVD err.\ (full)} & \textbf{Quant.} \\
\midrule
q\_proj & $2048{\times}2560$ & $1.1\!\times\!10^{-15}$ & 0.186 & $1.9\!\times\!10^{-15}$ & 0.0017 \\
k\_proj & $1024{\times}2560$ & $1.2\!\times\!10^{-15}$ & 0.249 & $1.9\!\times\!10^{-15}$ & 0.0019 \\
v\_proj & $1024{\times}2560$ & $8.1\!\times\!10^{-16}$ & 0.669 & $3.1\!\times\!10^{-15}$ & 0.0022 \\
o\_proj & $2560{\times}2048$ & $1.1\!\times\!10^{-15}$ & 0.311 & $2.4\!\times\!10^{-15}$ & 0.0020 \\
\bottomrule
\end{tabular}
\end{table}

\subsection{Rank Sweep}

\begin{table}[h]
\centering
\caption{SVD approximation error vs.\ rank for q\_proj. All chip-fidelity values
         are $\approx 10^{-15}$ (not shown; compilation is exact at all ranks).}
\label{tab:rank_sweep}
\begin{tabular}{rrrrrr}
\toprule
\textbf{Rank} & \textbf{Energy (\% Frob.)} & $\varepsilon_{\rm SVD}$ & \textbf{MZIs/mesh} &
\textbf{Area (std)} & \textbf{Area (cmp)} \\
\midrule
    64 & 28.8\% & 0.186 &        2,016 &    4.9\,mm$^2$ &    1.6\,mm$^2$ \\
   128 & 38.2\% & 0.166 &        8,128 &   19.7\,mm$^2$ &    6.6\,mm$^2$ \\
   256 & 52.8\% & 0.142 &       32,640 &   78.6\,mm$^2$ &   26.2\,mm$^2$ \\
   512 & 72.6\% & 0.108 &      130,816 &  314.6\,mm$^2$ &  104.9\,mm$^2$ \\
  1024 & 92.3\% & 0.061 &      523,776 & 1258.3\,mm$^2$ &  419.4\,mm$^2$ \\
  2048 & 100\%  & 0.000 &    2,096,128 & 5033.2\,mm$^2$ & 1677.7\,mm$^2$ \\
\bottomrule
\end{tabular}
\end{table}

\paragraph{Honest interpretation.}
The rank-64 SVD errors (0.19--0.67) are genuine information loss versus the
full weight matrix.
They do not indicate a compiler error.
The chip executes the rank-$r$ SVD approximation to machine precision.
End-to-end task accuracy (MedQA, PubMedQA) at each rank is a separate question
requiring full-model evaluation; this is not reported here.
The v\_proj matrix has the slowest singular value decay, requiring higher rank
for a given approximation quality --- this is a property of the weight matrix,
not of our compiler.

% ----------------------------------------------------------------
\section{Hardware Architecture}

\paragraph{Platform.} 220\,nm silicon-on-insulator (SOI), 1310\,nm wavelength,
TiN thermal phase shifters (10--20\,mW for $\pi$ shift), Ge photodetectors
($>$40\,GHz bandwidth).

\paragraph{Single chip.} 10\,mm $\times$ 10\,mm, 4,096 MZIs (rank-64 mesh in
8\,mm $\times$ 8\,mm), 64 fiber I/O channels per edge.

\paragraph{Power.}
$\approx 8$\,W total (laser + readout + FPGA) vs.\ $\approx 300$\,W for A100
GPU: \textbf{37$\times$ improvement}.
At 1 token/ms: 8\,mJ/token vs.\ 300\,mJ/token.

\paragraph{WDM parallelism.}
8 wavelengths at 1310\,nm $\pm$ 400\,GHz process all 8 attention heads
simultaneously on a single physical chip, reducing system chip count 8$\times$.

% ----------------------------------------------------------------
\section{Related Work}

Shen et al.~\citep{shen2017} first demonstrated deep learning with coherent
nanophotonic circuits on a small synthetic problem.
Bandyopadhyay et al.~\citep{bandyopadhyay2022} showed single-chip photonic
training.
Hamerly et al.~\citep{hamerly2019} analysed scalability and noise limits of
optical neural networks.
Pai et al.~\citep{pai2023} demonstrated 512-mode photonic processors with in-situ
calibration.
Our work differs in: (a) targeting a specific production medical AI model
(MedGemma-4B) rather than synthetic benchmarks; (b) the static compilation
approach that eliminates dynamic weight updates; (c) the full rank sweep from
64 to 2048 with honest error reporting; (d) an explicit deployment motivation
(KNH Nairobi, global-south hospitals).

% ----------------------------------------------------------------
\section{Conclusion}

PhotoMedGemma demonstrates that a production medical AI model (MedGemma-4B-IT)
can be compiled to photonic MZI phase settings with machine-precision fidelity.
The Clements algorithm faithfully encodes the target SVD unitaries at all
tested ranks (64--2048).
The rank-$r$ approximation error is a controllable design parameter, not a
compilation artefact.
The system projects 37$\times$ better energy efficiency than GPU inference,
enabling deployment at Kenyatta National Hospital and similar settings across
the global south.

Code, phase maps, sweep data, and this paper are at:
\url{github.com/GabuGravin41/photonmedgemma}.

% ----------------------------------------------------------------
\bibliographystyle{plain}
\begin{thebibliography}{9}

\bibitem{medgemma2024}
Google DeepMind. \textit{MedGemma: Medical AI built on Gemma.} arXiv:2507.05201, 2025.

\bibitem{shen2017}
Y.\ Shen et al. ``Deep learning with coherent nanophotonic circuits.''
\textit{Nature Photonics}, 11(7):441--446, 2017.

\bibitem{bandyopadhyay2022}
S.\ Bandyopadhyay et al. ``Single chip photonic deep neural network with
accelerated training.'' \textit{arXiv:2208.01623}, 2022.

\bibitem{clements2016}
W.R.\ Clements et al. ``Optimal design for universal multiport interferometers.''
\textit{Optica}, 3(12):1460--1465, 2016.

\bibitem{reck1994}
M.\ Reck et al. ``Experimental realization of any discrete unitary operator.''
\textit{Physical Review Letters}, 73(1):58, 1994.

\bibitem{hamerly2019}
R.\ Hamerly et al. ``Large-scale optical neural networks based on photoelectric
multiplication.'' \textit{Physical Review X}, 9(2):021032, 2019.

\bibitem{pai2023}
S.\ Pai et al. ``Experimentally realized in situ backpropagation for deep
learning in photonic neural networks.'' \textit{Science}, 380(6643):398--404, 2023.

\bibitem{gemma3_2024}
Google DeepMind. \textit{Gemma 3 Technical Report.} arXiv:2408.00118, 2024.

\end{thebibliography}

\end{document}
